\documentclass[12pt, a4paper]{article}

%=============== PAQUETES ===============
\usepackage[utf8]{inputenc}
\usepackage[T1]{fontenc}
\usepackage[spanish]{babel}
\usepackage{geometry}
\usepackage{amsmath}
\usepackage{graphicx}
\usepackage{xcolor}
\usepackage{hyperref}
\usepackage{tikz}
\usepackage{circuitikz}
\usetikzlibrary{arrows.meta, patterns, babel}

%=============== CONFIGURACIÓN ===============
\geometry{a4paper, margin=2.5cm}

\hypersetup{
    colorlinks=true,
    linkcolor=blue,
    urlcolor=cyan,
    pdftitle={Análisis de Examen de Electrónica},
}

\definecolor{darkgreen}{rgb}{0.0, 0.5, 0.0}
\definecolor{darkred}{rgb}{0.8, 0.0, 0.0}
\definecolor{darkorange}{rgb}{0.9, 0.5, 0}

\newcommand{\acierto}[1]{\textcolor{darkgreen}{\textbf{Acierto:} #1}}
\newcommand{\error}[1]{\textcolor{darkred}{\textbf{Error:} #1}}
\newcommand{\faltante}[1]{\textcolor{blue!70!black}{\textbf{Faltó añadir:} #1}}
\newcommand{\incompleto}[1]{\textcolor{darkorange}{\textbf{Incompleto:} #1}}

%=============== DOCUMENTO ===============

\begin{document}

\title{Análisis de Examen de Dispositivos Electrónicos}
\author{Evaluación de Arianna Isava}
\date{11 de Noviembre de 2025}
\maketitle

\section*{Análisis General}
La estudiante Arianna Isava demuestra una sólida comprensión de los conceptos teóricos fundamentales del diodo. Sus respuestas son, en general, precisas y bien fundamentadas. Muestra una capacidad de síntesis notable, aunque algunas preguntas no fueron respondidas o quedaron incompletas.

\hrulefill
\vspace{1cm}

\section{Pregunta 1: Diferencias entre diodo real e ideal}

\subsection*{Respuesta del Estudiante (Resumen)}
Describe el diodo ideal como un 'interruptor perfecto' con una gráfica en forma de 'L', y el diodo real como un dispositivo no lineal con una 'curva suave exponencial'.

\subsection*{Evaluación}
\begin{itemize}
    \item \acierto{La respuesta es excelente. La analogía del interruptor perfecto y la descripción de la gráfica en forma de 'L' para el diodo ideal son muy precisas y demuestran una clara comprensión conceptual.}
    \item \acierto{La descripción de la curva del diodo real como 'suave exponencial' también es correcta.}
    \item \faltante{Para una respuesta exhaustiva, podría haber mencionado la caída de tensión específica del diodo real (0.7V), la corriente de fuga en inversa y la tensión de ruptura, pero la respuesta proporcionada es de alta calidad.}
\end{itemize}

\begin{figure}[h!]
    \centering
    \begin{tikzpicture}[scale=0.8, xshift=-0.5cm]
        % Ejes
        \draw[->] (-3,0) -- (3,0) node[right] {$V_D$};
        \draw[->] (0,-1) -- (0,4) node[above] {$I_D$};
        
        % Diodo Ideal (rojo)
        \draw[thick, red] (0,0) -- (0,3.5);
        \draw[thick, red] (-2.5,0) -- (0,0);
        \node[red, below right] at (1,2) {Ideal};

        % Diodo Real (azul)
        \draw[thick, blue, domain=0:0.8, smooth] plot (\x, {0.05*(exp(8*\x) - 1)});
        \draw[thick, blue] (-2.5, -0.05) -- (0, -0.05);
        \node[blue, above right] at (1, 2.5) {Real};
        \draw[dashed] (0.7,0) -- (0.7, 1.5);
        \node[below] at (0.7,0) {0.7V};
    \end{tikzpicture}
    \caption{Comparación de las curvas I-V de un diodo real y uno ideal.}
\end{figure}

\newpage
\section{Pregunta 2: Funcionamiento de un circuito rectificador}

\subsection*{Respuesta del Estudiante (Resumen)}
El espacio para la respuesta fue dejado en blanco.

\subsection*{Evaluación}
\begin{itemize}
    \item \incompleto{La pregunta no fue respondida.}
\end{itemize}

\subsection*{Respuesta Correcta}
El rectificador de media onda simple consta de una fuente AC, un diodo y una resistencia de carga ($R_L$). Durante el semiciclo positivo de la señal de entrada, el diodo se polariza en directa y permite el paso de la corriente. Durante el semiciclo negativo, se polariza en inversa y bloquea la corriente, "recortando" esa parte de la señal en la salida.

\begin{figure}[h!]
    \centering
    \begin{circuitikz}[scale=0.9, transform shape]
        \draw (0,0) to[sV, l=$V_{in}$] (0,2) to[short] (1,2) to[D, l=D] (3,2) to[R, l=$R_L$] (5,2) to[short] (6,2) node[ground]{};
        \draw (0,0) to[short] (6,0) to[short] (6,2);
    \end{circuitikz}
    \caption{Diagrama de un rectificador de media onda simple.}
\end{figure}

\begin{figure}[h!]
    \centering
    \begin{tikzpicture}[scale=1.2]
        % Señal de entrada
        \draw[->] (-0.2,0) -- (4,0) node[below] {$t$};
        \draw[->] (0,-1.2) -- (0,1.2) node[left] {$V$};
        \draw[blue, thick] (0,0) sin (1,1) cos (2,0) sin (3,-1) cos (4,0);
        \node[above] at (2,1.2) {Entrada ($V_{in}$)};

        % Señal de salida
        \begin{scope}[yshift=-2.5cm]
            \draw[->] (-0.2,0) -- (4,0) node[below] {$t$};
            \draw[->] (0,-1.2) -- (0,1.2) node[left] {$V_{out}$};
            \draw[red, thick] (0,0) sin (1,1) cos (2,0);
            \draw[red, thick, dashed] (2,0) sin (3,-1) cos (4,0); % Semiciclo negativo recortado
            \draw[red, thick] (2,0) -- (4,0);
            \node[above] at (2,1.2) {Salida ($V_{out}$)};
        \end{scope}
    \end{tikzpicture}
    \caption{Formas de onda de entrada y salida para un rectificador de media onda.}
\end{figure}

\section{Pregunta 3: Relación corriente-tensión en un diodo}

\subsection*{Respuesta del Estudiante (Resumen)}
Presenta correctamente la ecuación de Shockley ($I_D = I_S(e^{V_D/V_T} - 1)$) y explica que para polarización directa, el término exponencial domina, resultando en la aproximación $I_D \approx I_S e^{V_D/V_T}$.

\subsection*{Evaluación}
\begin{itemize}
    \item \acierto{La ecuación fundamental del diodo está escrita correctamente.}
    \item \acierto{La explicación de la aproximación para polarización directa es conceptualmente correcta y demuestra entendimiento del comportamiento de la ecuación.}
    \item \faltante{Podría haber explicado el significado de cada término en la ecuación ($I_S$: corriente de saturación inversa, $V_T$: tensión térmica) y mencionar brevemente el comportamiento en polarización inversa, donde $I_D \approx -I_S$.}
\end{itemize}

\section{Pregunta 4: Ventaja del rectificador de onda completa}

\subsection*{Respuesta del Estudiante (Resumen)}
Indica que el rectificador de onda completa es preferible porque permite una 'salida de corriente más estable y eficiente', mientras que el de media onda se usa en aplicaciones de baja potencia por su 'voltaje de salida bajo y poca regulación'.

\subsection*{Evaluación}
\begin{itemize}
    \item \acierto{La respuesta es muy buena. Identifica correctamente la mayor estabilidad y eficiencia como las ventajas clave del rectificador de onda completa.}
    \item \acierto{La caracterización del rectificador de media onda como adecuado para aplicaciones simples y de baja potencia también es correcta.}
    \item \faltante{Para mayor rigor técnico, podría haber explicado que la 'mayor estabilidad' se debe a que la frecuencia del rizado es el doble, lo que facilita el filtrado.}
\end{itemize}

\section{Pregunta 5: Modelo de tramos lineales}

\subsection*{Respuesta del Estudiante (Resumen)}
Describe el modelo como una 'aproximación lineal por partes de la curva I-V real', que representa el comportamiento del diodo con 'segmentos lineales con ecuaciones simples' usando un 'modelo de dos tramos'.

\subsection*{Evaluación}
\begin{itemize}
    \item \acierto{La respuesta es conceptualmente correcta. Entiende que es una aproximación por segmentos de la curva real para simplificar el análisis.}
    \item \incompleto{La descripción es algo genérica. No especifica en qué consisten esos 'dos tramos' (es decir, un interruptor abierto por debajo de $V_{th}$ y una fuente de tensión en serie con una resistencia por encima de $V_{th}$).}
    \item \error{La frase final 'describir el comportamiento del voltaje del diodo en el tiempo' es confusa. El modelo describe la relación estática I-V, no su evolución temporal, aunque se use para analizar circuitos que varían en el tiempo.}
\end{itemize}

\begin{figure}[h!]
    \centering
    \begin{tikzpicture}[scale=0.8]
        % Ejes
        \draw[->] (-1,0) -- (3,0) node[right] {$V_D$};
        \draw[->] (0,-0.5) -- (0,4) node[above] {$I_D$};
        
        % Modelo lineal por tramos
        \draw[thick, blue] (0.7,0) -- (2, 2.6); % Resistencia r_D > 0
        \draw[thick, blue] (-0.5,0) -- (0.7,0);
        \node[blue, right] at (1.5, 3) {Modelo Lineal por Tramos};
        \draw[dashed] (0.7,0) -- (0.7, 2);
        \node[below] at (0.7,0) {$V_{th}$};
    \end{tikzpicture}
    \caption{Característica I-V del modelo de tramos lineales.}
\end{figure}

\end{document}